% chapters/ch08-naming-documentation.tex
\chapter{Naming, Documentation, and Design Hygiene}

\begin{learningobjectives}
In this chapter you will learn how naming and documentation affect correctness.
You will be able to choose consistent names, document schemas at the right level of detail, and write small design records that make your modelling decisions auditable.
\end{learningobjectives}

\section{Why hygiene matters}
A relational schema is a long-lived artifact.
Even small projects accumulate time: new features, new people, new interpretations.
When a schema is unclear, the database becomes a place where meaning drifts silently.

This is why naming and documentation are not cosmetic.
They are part of correctness.
If two developers interpret the same column differently, the system will store contradictory facts while still “working”.
Good hygiene reduces interpretability, and therefore reduces hidden bugs.

\begin{definitionbox}
\term{Design hygiene} is the practice of keeping naming, structure, and documentation consistent so that the schema has a single stable interpretation.
\end{definitionbox}

\section{Naming principles}
Naming is about reducing ambiguity.
A reader should guess the meaning of a relation or attribute correctly without needing a meeting.

\begin{principlebox}
\textbf{Naming goals}\\[2pt]
Names should be:
\begin{itemize}[leftmargin=*,nosep]
  \item consistent across the schema,
  \item specific enough to avoid multiple interpretations,
  \item stable over time (avoid names tied to temporary implementation choices).
\end{itemize}
\end{principlebox}

\subsection{Entity names}
Use singular nouns for entity relations: \rel{Patient}, \rel{Doctor}, \rel{Appointment}, \rel{Loan}.
This matches the idea “a row represents one instance”.

For association relations, use compound names that communicate meaning:
\rel{DoctorSpecialty}, \rel{StudentCourseEnrollment}, \rel{OrderLine}.
Avoid generic names like \rel{Mapping} or \rel{Link} because they hide semantics.

\subsection{Attribute names}
Attribute names should communicate:
what the value is, and what it refers to.

A reliable convention is:
\begin{itemize}[leftmargin=*,nosep]
  \item primary key: \texttt{<entity>\_id} (e.g., \texttt{patient\_id}),
  \item foreign key: same naming pattern, matching the referenced entity,
  \item booleans: \texttt{is\_*}, \texttt{has\_*} (e.g., \texttt{is\_active}),
  \item timestamps/dates: include the semantic moment.\par
        Example: \texttt{checkout\_date}, \texttt{return\_date}, \texttt{created\_at}.
\end{itemize}

Avoid names like \texttt{date} or \texttt{status} when multiple “dates” or “statuses” exist in the same domain.
Ambiguous names force readers to search for meaning and increase the chance of misuse.

\subsection{Do not encode meaning in magic abbreviations}
Abbreviations can be useful, but they also hide meaning.
If you use abbreviations, keep them standard and document them.
Avoid project-specific abbreviations in the schema unless the domain itself uses them.

\section{Naming conventions that scale}
Small naming decisions become expensive when a schema grows.
Use conventions that survive more tables, more teams, and more time.

\subsection{Identifiers}
Use \texttt{<entity>\_id} for primary keys and match foreign keys to the referenced entity.
This keeps joins readable and reduces misinterpretation.

\subsection{Timestamps and dates}
Name dates by their semantic moment: \texttt{created\_at}, \texttt{checkout\_date}, \texttt{valid\_from}.
Avoid generic \texttt{date} or \texttt{timestamp} in domains with multiple time concepts.

\subsection{Booleans and state fields}
Use \texttt{is\_*} or \texttt{has\_*} for booleans.
For multi-state attributes, use a noun like \texttt{status} with a documented domain.
If there are multiple statuses, name them explicitly (e.g., \texttt{loan\_status}, \texttt{membership\_status}).

\section{Documentation layers}
Documentation fails in two ways.
It can be too small, leaving meaning undefined.
Or it can be too large, becoming outdated.

Use layers so that each artifact has a clear purpose:
\begin{itemize}[leftmargin=*,nosep]
  \item README-level overview: domain narrative and goals.
  \item Schema specification: relations, keys, and constraints.
  \item Decision records: why modelling choices were made.
\end{itemize}

\begin{figure}[htbp]
\centering
\includegraphics[width=0.8\textwidth]{documentation-layers.pdf}
\caption{Layered documentation keeps purpose, structure, and decisions aligned.}
\end{figure}

This set is stable because it explains \emph{meaning}, not implementation details.

\begin{notebox}
Avoid documenting things that will change frequently unless they are generated automatically.
For example, detailed index plans or performance notes often drift.
Meaning and constraints drift less, and are more valuable in early design.
\end{notebox}

\section{Schema spec block you can paste into any project}
A compact schema block prevents most confusion.
It forces you to state purpose, identity, and constraints in a repeatable format.

\begin{examplebox}
\textbf{Schema spec block}\\[4pt]
Relation: \rel{Appointment}\\
Purpose: Stores appointments scheduled between a patient and a doctor.\\
Primary key: \texttt{appointment\_id} (surrogate).\\
Foreign keys: \texttt{patient\_id} \(\rightarrow\) \rel{Patient} (mandatory), \texttt{doctor\_id} \(\rightarrow\) \rel{Doctor} (mandatory).\\
Key constraints: start\_time mandatory; no double booking per doctor at same start\_time (slot-based).\\
Notes: If appointment duration varies, double booking becomes an interval overlap constraint.
\end{examplebox}

\begin{examplebox}
\textbf{Constraint block template}\\[4pt]
C1: Appointment.patient\_id references an existing Patient and is mandatory.\\
C2: Appointment.doctor\_id references an existing Doctor and is mandatory.\\
C3: Appointment.start\_time is mandatory.\\
C4: No two appointments may share the same (doctor\_id, start\_time) if time is slot-based.
\end{examplebox}

When a constraint cannot be enforced directly, document how violations are detected.
This keeps correctness visible even when enforcement is external.

\section{Decision records: making modelling choices auditable}
Some decisions are not obvious from the final schema.
A decision record is a short note that explains why a choice was made.
It is especially useful for identity choices, normalization boundaries, and state modelling.

A decision record should include:
\begin{itemize}[leftmargin=*,nosep]
  \item the decision,
  \item the alternatives considered,
  \item the rationale,
  \item and the implications.
\end{itemize}

\begin{examplebox}
\textbf{Decision record: Title vs Copy}\\[4pt]
Decision: Model \rel{Title} and \rel{Copy} as separate entities.\\
Alternatives: Single relation with one row per physical copy containing title metadata.\\
Rationale: Avoid duplication of title facts across copies; support availability and loan history; support multiple copies per title.\\
Implications: Loans reference copies; reservations reference titles; requires an “active loan” constraint per copy.
\end{examplebox}

\begin{examplebox}
\textbf{Decision record: Surrogate vs natural key}\\[4pt]
Decision: Use a surrogate key for \rel{Patient}.\\
Alternatives: Use national health number as primary key.\\
Rationale: Identifiers may be missing at first contact; surrogate keys keep identity stable across merges.\\
Implications: National health number stored as optional unique attribute; matching rules documented separately.
\end{examplebox}

This format is useful in teaching because it shows that design is deliberate.
It also mirrors how real teams justify schema decisions.

\section{Common mistakes}
Naming and documentation mistakes are easy to introduce.

Mistake: Using ambiguous attribute names (e.g., \texttt{date}, \texttt{status}).\\
Fix: Name attributes by their semantic role (\texttt{checkout\_date}, \texttt{loan\_status}).

Mistake: Mixing naming styles across the schema.\\
Fix: Choose one convention for identifiers and use it everywhere.

Mistake: Writing decision records only after problems appear.\\
Fix: Capture key decisions when you make them, even if brief.

Mistake: Documenting implementation details instead of meaning.\\
Fix: Prioritize rules, constraints, and identity.

\section{Mini checklist}
Before you share a schema, check:
\begin{itemize}[leftmargin=*,nosep]
  \item Are relation names singular and meaningful?
  \item Do identifiers follow a consistent pattern (\texttt{<entity>\_id})?
  \item Are date/time attributes named by their semantic moment?
  \item Are boolean attributes named with \texttt{is\_*} or \texttt{has\_*}?
  \item Is each multi-state field documented with allowed values?
  \item Are key decisions recorded in short decision records?
  \item Do schema blocks include purpose, keys, and constraints?
  \item Is documentation focused on meaning rather than implementation?
\end{itemize}

\begin{chaptersummary}
Design hygiene protects meaning.
Consistent naming reduces ambiguity, and compact documentation makes modelling decisions stable over time.
A good schema specification includes rules, constraints, and traceability, not just relation names.
Decision records make the non-obvious choices auditable and teach readers how to reason about trade-offs.
\end{chaptersummary}

\begin{exercisebox}
Answer in full sentences.
\begin{enumerate}[label=\arabic*.]
  \item Choose a domain and propose naming conventions for relations and identifiers. Give three examples.
  \item Write relation-level documentation blocks for two relations in your schema.
  \item Write a constraint block (C1, C2, ...) for one relation with at least five constraints.
  \item Write one decision record for a modelling decision that could be debated (e.g., surrogate vs natural key, denormalization boundary, state modelling).
  \item Review your schema and rename at least three attributes to reduce ambiguity. Explain why.
  \item Create a mini schema spec block for one relation in your project using the template above.
  \item Identify two ambiguous names in a domain you know and propose better alternatives.
  \item Write a documentation layer outline (README, schema spec, decision records) for a small project.
\end{enumerate}
\end{exercisebox}

In the next chapter we apply these practices in a clinic case study.
Naming and documentation choices will be part of the model, not an afterthought.
